%%%%%%%%%%%%%%%%%%%%%%%%%%%%%%%%%%%%%%%%%%%%%%%%%%%%%%%%%%%%%%%%%%%%%%%%%%%
%%  chapter5.tex  –  Глава 5: Программная платформа RobustIDPS.ai
%%  v9 – Глава промышленного развёртывания: целевые пользователи, 46-стр. PWA,
%%       9 функциональных групп, 13+ моделей, защита LLM, PQC-IDS,
%%       плагин Snort/Suricata, 8 рисунков для скриншотов.
%%       Только заголовки \section добавлены в оглавление.
%%%%%%%%%%%%%%%%%%%%%%%%%%%%%%%%%%%%%%%%%%%%%%%%%%%%%%%%%%%%%%%%%%%%%%%%%%%

\chapter{RobustIDPS.ai: Программная платформа для промышленного развёртывания}
\label{ch:platform}

В настоящей главе описана программная платформа с открытым исходным кодом RobustIDPS.ai~\cite{Anaedevha2026robustidps} (DOI: 10.5281/zenodo.19129512; исходный код: \texttt{https://github.com/rogerpanel/robustidps.ai}), реализующая все методы, разработанные в Главе~\ref{ch:methods}, как единую развёртываемую систему. В главе определены целевые пользователи платформы, описаны её архитектура и функциональные группы, подробно представлены 13+ интегрированных моделей нейронных сетей, представлены возможности SOC-аналитики и тестирования безопасности LLM, описан модуль Multi-Agent PQC-IDS, продемонстрирована интеграция с существующими системами СОПВ (Snort, Suricata) и обобщены достижения платформы. Не менее шести скриншотов пользовательского интерфейса включены для иллюстрации операционных возможностей.


%%=====================================================================
\section{Обзор платформы и целевые пользователи}
\label{sec:platform_overview}
%%=====================================================================

RobustIDPS.ai --- комплексная ИИ-нативная платформа обнаружения и предотвращения вторжений, содержащая 46 взаимосвязанных страниц, организованных в 9 функциональных групп, интегрирующая более 13 моделей машинного и глубокого обучения для классификации сетевого трафика по 34 классам атак и 83 инженерным признакам. Платформа решает три ограничения существующих коммерческих аналогов (CrowdStrike Charlotte AI, SentinelOne Purple AI, Microsoft Security Copilot, Google Sec-Gemini~\cite{Smith2023}): проприетарная интеграция LLM одного вендора, препятствующая межпровайдерному бенчмаркингу, отсутствие систематической оценки состязательной робастности и отсутствие поддержки анализа постквантового криптографического трафика.

Платформа предназначена для пяти категорий целевых пользователей:

\noindent\textbf{SOC-аналитики (Уровни 1--3).} Персонал центра операций безопасности, осуществляющий мониторинг, расследование и реагирование на инциденты безопасности. Платформа обеспечивает агентное авторасследование, поиск угроз на естественном языке, автоматическую генерацию отчётов об инцидентах и сортировку оповещений на основе неопределённости, снижающую среднее время сортировки на 34\%.

\noindent\textbf{Инженеры кибербезопасности.} Специалисты по сетевой безопасности, развёртывающие и настраивающие системы обнаружения вторжений. Платформа генерирует развёртываемые правила Snort/Suricata из обнаруженных угроз и обеспечивает мониторинг в реальном времени через потоковую передачу WebSocket.

\noindent\textbf{Исследователи MLSecOps.} Исследователи безопасности машинного обучения, оценивающие робастность ИИ-компонентов. Платформа обеспечивает интегрированное состязательное тестирование (6 семейств атак), оценку поверхности атак LLM (5 модулей) и панели соответствия MLSecOps (MITRE ATT\&CK, OWASP LLM Top 10, NIST AI RMF).

\noindent\textbf{Разработчики ИИ.} Исследователи и инженеры, создающие модели обнаружения на основе графов. Платформа предоставляет 13+ предобученных моделей с настраиваемыми гиперпараметрами, студию абляции и воспроизводимые бенчмарк-наборы данных.

\noindent\textbf{Академические исследователи.} Аспиранты, постдокторанты и преподаватели, проводящие исследования в области робастности графовых нейронных сетей. Платформа служит воспроизводимым исследовательским инструментом с DOI-архивированным кодом и наборами данных.


%%=====================================================================
\section{Архитектура системы}
\label{sec:system_arch}
%%=====================================================================

RobustIDPS.ai развёрнута как прогрессивное веб-приложение (PWA) со следующим технологическим стеком:

\noindent\textbf{Фронтенд.} Одностраничное приложение React 18 с компонентами Material-UI, адаптивным дизайном и манифестом PWA для развёртывания с поддержкой офлайн-режима через кэширование сервис-воркером.

\noindent\textbf{Бэкенд.} Бэкенд Flask/Python 3.11+ с 60+ RESTful API эндпоинтами, постоянными WebSocket-каналами для мониторинга в реальном времени и Celery для асинхронного выполнения задач (обучение моделей, пакетный анализ).

\noindent\textbf{Слой данных.} PostgreSQL 16 для реляционного хранения (профили пользователей, журналы обнаружений, записи расследований), Redis 7 для кэширования и управления сессиями, файловое хранилище для загруженных PCAP/CSV-наборов данных.

\noindent\textbf{Оркестрация.} Docker Compose обеспечивает развёртывание одной командой с конфигурацией на основе переменных окружения для профилей разработки, промежуточного этапа и производства.

\begin{figure}[H]
\centering
\includegraphics[width=0.92\textwidth]{figures/ridps_arch_v10.pdf}
\caption{Высокоуровневая архитектура системы RobustIDPS.ai, показывающая все 9 функциональных групп, поддерживающую инфраструктуру (Docker Compose, PostgreSQL, Redis, WebSocket, Celery) и поток данных от приёма сырого трафика через обнаружение, анализ и реагирование.}
\label{fig:platform_arch}
\end{figure}


%%=====================================================================
\section{Функциональные группы и реестр страниц}
\label{sec:functional_groups}
%%=====================================================================

46 страниц платформы организованы в 9 функциональных групп, обобщённых в Таблице~\ref{tab:functional_groups}.

\begin{table}[htbp]
\centering
\caption{Функциональные группы RobustIDPS.ai (46 страниц в 9 группах).}
\label{tab:functional_groups}
\small
\begin{tabular}{lcp{7.5cm}}
\toprule
\textbf{Функциональная группа} & \textbf{Стр.} & \textbf{Ключевые возможности} \\
\midrule
Командный центр ИИ      & 5 & Панель, Загрузка и анализ, Аналитика, Живой монитор, Каузальность оповещений \\
Активная защита ИИ      & 5 & Арена красной команды, Реагирование на угрозы, SOC-копилот, Федеративное обучение, Предиктор цепочек атак \\
SOC-аналитика           & 6 & Авторасследование, Поиск угроз, Отчёты об инцидентах, Разведка угроз, Генератор правил, Маппер CVE \\
Поверхности атак LLM    & 5 & Инъекция промптов, Таксономия джейлбрейков, Отравление RAG, Мультиагентная цепочка, Отравление данных \\
Стандарты MLSecOps       & 3 & Маппер MITRE ATT\&CK, Сортировка оповещений, Хаб соответствия \\
Безопасность и управл. ИИ & 4 & PQ-криптография, Zero-Trust, Цепочка поставок, Исследовательский хаб \\
Новые методы ИИ         & 5 & Непрерывное обучение, RL-агент реагирования, Состязат. оценка, Объяснимость, Автокодировщик \\
Данные и модели ИИ      & 3 & Наборы данных (6), Модели (13+), Студия абляции \\
Система                 & 5 & Профиль/RobustID, Администрирование, Архитектура, Документация API, О платформе \\
\midrule
\textbf{Всего}          & \textbf{46} & \\
\bottomrule
\end{tabular}
\end{table}


%%=====================================================================
\section{Командный центр ИИ}
\label{sec:command_centre}
%%=====================================================================

Командный центр ИИ --- основной операционный интерфейс, включающий пять страниц для сквозного анализа трафика от загрузки до разрешения оповещений.

\subsection{Панель управления}

Основная панель отображает статистику обнаружения в реальном времени: общее число проанализированных потоков, распределение атак по категориям, метрики точности моделей и разбивку оповещений по степени серьёзности. Интерактивные диаграммы временных рядов показывают тенденции обнаружения по настраиваемым окнам (от 1 часа до 30 дней).

\begin{figure}[H]
\centering
\includegraphics[width=0.96\textwidth]{figures/Dashboard_v10.png}
\caption{Панель управления RobustIDPS.ai, отображающая статистику обнаружения в реальном времени, распределение категорий атак и метрики производительности моделей. Левая панель показывает активную конфигурацию модели; центр --- временную шкалу обнаружений; правая панель --- разбивку оповещений по степени серьёзности.}
\label{fig:screenshot_dashboard}
\end{figure}

\subsection{Загрузка и анализ}

Пользователи загружают PCAP или CSV файлы через интерфейс перетаскивания. Модуль \texttt{FeatureExtractor} вычисляет 83-мерный вектор признаков из сырых пакетных данных, применяет z-стандартизацию и направляет признаки в выбранную модель обнаружения. Результаты представлены с попоточной классификацией, оценками уверенности, декомпозицией неопределённости и сопоставлениями техник MITRE ATT\&CK.

\begin{figure}[H]
\centering
\includegraphics[width=0.92\textwidth]{figures/Dashboard_v11.png}
\caption{Интерфейс загрузки и анализа: (вверху) панель загрузки файлов с выпадающим списком выбора модели, (центр) таблица результатов обнаружения с попоточной классификацией и оценками уверенности, (внизу) разбивка по категориям атак с метриками производительности по классам.}
\label{fig:screenshot_upload}
\end{figure}

\subsection{Живой монитор}

Потоковая передача в реальном времени на основе WebSocket отображает решения классификации по мере поступления графовых событий с настраиваемым окном удержания и субсекундными push-уведомлениями для оповещений высокой серьёзности. Живой монитор поддерживает одновременное отображение выходов нескольких моделей для ансамблевого сравнения.

\subsection{Каузальность оповещений}

Страница каузальности оповещений строит графы темпоральных зависимостей между связанными оповещениями, позволяя аналитикам прослеживать многоэтапные кампании атак. Корреляция основана на общих IP источника/назначения, временной близости ($\pm$30\,с) и связках техник MITRE ATT\&CK.


%%=====================================================================
\section{Модели обнаружения}
\label{sec:platform_models}
%%=====================================================================

Платформа интегрирует более 13 архитектур нейронных сетей, реализующих все семь методов из Главы~\ref{ch:methods}. Таблица~\ref{tab:platform_models} обобщает их производительность на валидационном бенчмарке CIC-IoT-2023.

\begin{table}[htbp]
\centering
\caption{Производительность моделей обнаружения на валидационном бенчмарке CIC-IoT-2023.}
\label{tab:platform_models}
\small
\begin{tabular}{lcccc}
\toprule
\textbf{Модель} & \textbf{Точность} & \textbf{F1} & \textbf{ECE} & \textbf{Задержка} \\
\midrule
SurrogateIDS (7-Branch)   & 97.8\% & 0.976 & 0.031 & 0.8\,мс \\
CL-RL Unified + EWC      & 96.8\% & 0.965 & 0.035 & 1.1\,мс \\
CT-TGNN~\cite{Anaedevha2025cais} (М1)              & 96.2\% & 0.958 & 0.049 & 1.2\,мс \\
SDE-TGNN~\cite{Anaedevha2026sdetgnn}                  & 96.0\% & 0.955 & 0.038 & 1.4\,мс \\
MambaShield~\cite{Anaedevha2026mamba} (М4)          & 95.9\% & 0.954 & 0.042 & 0.5\,мс \\
Stochastic Transformer (М5) & 95.4\% & 0.949 & 0.078 & 1.8\,мс \\
FedGTD~\cite{Anaedevha2026neurips} (М2)               & 95.1\% & 0.946 & 0.040 & 1.3\,мс \\
CyberSecLLM~\cite{Anaedevha2026cybersecllm} \CyberSecLLM               & 97.1 97.1\% & 0.965 & 0.025 & 6.8\,мс \\
Multi-Agent PQC-IDS       & 94.3\% & 0.938 & 0.044 & 1.5\,мс \\
Neural ODE                & 94.8\% & 0.938 & 0.049 & 8.7\,мс \\
Optimal Transport         & 93.9\% & 0.930 & 0.052 & 3.4\,мс \\
Variational Autoencoder   & 92.4\% & 0.914 & 0.061 & 2.1\,мс \\
Adversarial Autoencoder   & 91.8\% & 0.908 & 0.058 & 2.3\,мс \\
\bottomrule
\end{tabular}
\end{table}

MambaShield достигает наименьшей задержки на поток (0,5\,мс) благодаря формулировке пространства состояний с линейной сложностью, что делает его предпочтительным для высокопроизводительных развёртываний. CyberSecLLM достигает наименьшей ECE (0,025) благодаря калибровочным свойствам вероятностных выходов языковой модели. Ансамбль SurrogateIDS обеспечивает лучший компромисс точность--калибровка и служит магистральным компонентом обнаружения по умолчанию.

\begin{figure}[H]
\centering
\includegraphics[width=0.92\textwidth]{figures/Models_v10.png}
\caption{Страница моделей, отображающая 13+ интегрированных моделей обнаружения с настраиваемыми гиперпараметрами, метриками точности по каждой модели и бенчмарками задержки. Пользователи могут выбрать любую модель в качестве активного детектора или запустить ансамблевые сравнения.}
\label{fig:screenshot_models}
\end{figure}


%%=====================================================================
\section{Пакет SOC-аналитики}
\label{sec:soc_intelligence}
%%=====================================================================

Пакет SOC-аналитики включает шесть целевых страниц, преобразующих сырой выход обнаружения в действенную аналитику для аналитиков уровней~1--3.

\subsection{Агентное авторасследование}

Конвейер выполняет четыре фазы: (1)~Сбор свидетельств --- извлечение инициирующего потока, $\pm$30\,с темпорального контекста и коррелированных оповещений; (2)~Контекстуальное обогащение --- запросы репутации IP, геолокации и частоты оповещений; (3)~Рассуждение LLM --- передача обогащённых свидетельств выбранной LLM для анализа первопричин и сопоставления MITRE ATT\MITRE ATT\&CK mappingCK~\cite{Hernandez2015} mapping; (4)~Синтез вердикта --- агрегация выхода в форматированный отчёт расследования.

\subsection{Поиск угроз на естественном языке}

Аналитики формулируют запросы в свободной форме (например, <<высокоуверенная DDoS на порту 443 с внешних IP>>). Парсер извлекает шесть измерений фильтрации (attack\_type, confidence, port, severity, protocol, destination) и применяет их к журналу обнаружений.

\subsection{Генератор правил СОВ}

Транслирует обнаруженные атаки в развёртываемые правила Suricata/Snort с использованием 21 шаблона, покрывающего разведку, DoS, эксплуатацию, латеральное перемещение и эксфильтрацию. Выход включает SID, приоритет и ссылки CVE.

\begin{figure}[H]
\centering
\includegraphics[width=0.96\textwidth]{figures/SOC_v10.png}
\caption{Пакет SOC-аналитики: (слева) панель авторасследования с выходом рассуждения LLM и сопоставлением MITRE ATT\MITRE ATT\&CK mappingCK~\cite{Hernandez2015} mapping, (центр) интерфейс поиска угроз с запросом на естественном языке, (справа) сгенерированное правило Suricata с SID, приоритетом и ссылками CVE.}
\label{fig:screenshot_soc}
\end{figure}


%%=====================================================================
\section{Тестирование безопасности LLM}
\label{sec:llm_testing}
%%=====================================================================

Модуль поверхностей атак LLM предоставляет пять интерактивных тестовых сред для оценки робастности рабочих процессов безопасности, интегрированных с LLM. Четырёхуровневый DefensePipeline~\cite{Anaedevha2026neurips,Brown2020} достигает общего обнаружения 94,5\% по 517 сценариям атак с четырьмя коммерческими бэкендами LLM (Claude 96,1\%, GPT-4o 94,8\%, Gemini 93,2\%, DeepSeek 91,5\%).

\begin{table}[htbp]
\centering
\caption{Показатели обнаружения атак LLM по модулям.}
\label{tab:llm_detection}
\small
\begin{tabular}{lcc}
\toprule
\textbf{Модуль атак} & \textbf{Категории} & \textbf{Обнаруж. \%} \\
\midrule
Инъекция промптов       & 8  & 96,3 \\
Таксономия джейлбрейков & 8  & 95,1 \\
Отравление RAG          & 4  & 93,7 \\
Мультиагентная цепочка  & 5  & 91,8 \\
Отравление данных       & 4  & 95,6 \\
\midrule
\textbf{Всего}          & \textbf{29} & \textbf{94,5} \\
\bottomrule
\end{tabular}
\end{table}

Пять модулей охватывают: площадку инъекции промптов (8 категорий инъекций, тестируемых против работающего DefensePipeline), таксономию джейлбрейков (8 техник, включая DAN-стиль переопределения, обход base64, контрабанду токенов), симулятор отравления RAG (4 вектора: инъекция документов, возмущение вложений, манипуляция метаданными, эксплуатация границ фрагментов), симулятор мультиагентных цепочечных атак (5 паттернов на 4-агентной топологии) и симулятор отравления данных (4 стратегии: переворот меток, инъекция бэкдоров, градиентные атаки, чистометочные атаки).

\begin{figure}[H]
\centering
\includegraphics[width=0.96\textwidth]{figures/LLM_v10.png}
\caption{Интерфейс тестирования безопасности LLM, показывающий площадку инъекции промптов с оценкой DefensePipeline в реальном времени. Левая панель показывает инъекционный вход; центр отображает санитизированный выход с флагами обнаружения и сопоставленными правилами; правая панель показывает статистику обнаружения по всем 4 бэкендам LLM.}
\label{fig:screenshot_llm}
\end{figure}


%%=====================================================================
\section{Соответствие стандартам MLSecOps}
\label{sec:mlsecops}
%%=====================================================================

Платформа интегрирует три основные системы безопасности и управления ИИ: сопоставление MITRE ATT\MITRE ATT\&CK mappingCK~\cite{Hernandez2015} mapping (30 классов обнаружения, сопоставленных техникам цепочки поражения с интерактивной визуализацией), таблицу оценки OWASP LLM Top~10~\cite{Xiong2020} (соответствие изданию 2025 года с оценками остаточного риска) и панель управления NIST AI Risk Management Framework (функции Govern, Map, Measure, Manage с свидетельствами по подкатегориям). ML-ассистированный классификатор сортировки оповещений назначает уровни серьёзности (Критический/Высокий/Средний/Низкий/Информационный) на основе класса атаки, уверенности, критичности актива и исторических показателей ложных срабатываний.


%%=====================================================================
\section{Модуль Multi-Agent PQC-IDS}
\label{sec:pqc_module}
%%=====================================================================

Модуль мультиагентной постквантовой криптографической СОВ~\cite{Anaedevha2026mapqc,NIST2024} представляет возможность прямо совместимого обнаружения Метода М4 (MambaShield), расширенного на среды постквантовых протоколов. Четыре специализированных агента взаимодействуют через слияние на основе взвешенного внимания:

\noindent\textbf{Аналитик трафика} --- извлекает признаки уровня потоков из трафика PQC-рукопожатий (размеры обмена ключами, паттерны согласования шифров, длины цепочек сертификатов).

\noindent\textbf{Специалист PQC} --- классифицирует семейства криптографических алгоритмов (CRYSTALS-Kyber, CRYSTALS-Dilithium, SPHINCS+, BIKE) и обнаруживает аномальные выборы параметров.

\noindent\textbf{Детектор аномалий} --- применяет анализ статистических отклонений тайминга потоков, энтропии полезной нагрузки и метаданных сессий.

\noindent\textbf{Координатор} --- объединяет выходы агентов через обученные веса внимания и формирует итоговую классификацию с оценками вклада каждого агента.

Полная модель содержит 70\,598 обучаемых параметров. Наборы данных PQC-трафика получены из Kaggle (DOI: 10.34740/kaggle/dsv/15424420) и опубликованы отдельно по адресу \texttt{github.com/rogerpanel/Multi-Agent-PQC-models}.


%%=====================================================================
\section{Развёртывание и интеграция с существующими СОПВ}
\label{sec:deployment}
%%=====================================================================

RobustIDPS.ai предназначена для развёртывания как плагин для существующих систем обнаружения вторжений, в частности Snort~3~\cite{Roesch1999} и Suricata. Интеграция работает через два механизма:

\noindent\textbf{Генерация правил.} Генератор правил СОВ транслирует обнаруженные угрозы в развёртываемые правила Suricata/Snort с использованием 21 шаблона, позволяя ML-обнаружению формировать действенные выходы в форматах, которые существующая инфраструктура уже потребляет.

\noindent\textbf{Потребление потока оповещений.} Платформа потребляет журналы оповещений Suricata/Snort (формат EVE JSON или unified2) как дополнительный вход, коррелируя сигнатурные обнаружения с аномальными классификациями ГНС для снижения ложных срабатываний и контекстуального обогащения.

Оркестрация Docker Compose обеспечивает развёртывание одной командой. Переключатель \texttt{fast\_mode} снижает задержку вывода на 60\% (отключая MC-дропаут и оценку неопределённости), достигая пропускной способности свыше 12\,000 потоков/секунду на одном ядре CPU --- достаточно для анализа каналов 10\,Гбит/с в реальном времени при соответствующей параллелизации.


%%=====================================================================
\section{Сравнение с коммерческими и открытыми платформами}
\label{sec:platform_comparison}
%%=====================================================================

\begin{table}[htbp]
\centering
\caption{Сравнение возможностей с открытыми и коммерческими платформами.}
\label{tab:platform_comparison}
\footnotesize
\begin{tabular}{p{3.5cm}ccccc}
\toprule
\textbf{Возможность} & \rotatebox{70}{\textbf{RobustIDPS.ai}} & \rotatebox{70}{\textbf{Suricata}} & \rotatebox{70}{\textbf{Snort 3}} & \rotatebox{70}{\textbf{CrowdStrike}} & \rotatebox{70}{\textbf{MS Copilot}} \\
\midrule
ML-обнаружение (13+ мод.) & \checkmark & ---        & ---        & Частично   & ---        \\
Состязат. робастность      & \checkmark & ---        & ---        & ---        & ---        \\
Непрерывное обучение       & \checkmark & ---        & ---        & ---        & ---        \\
Мультипровайд. LLM         & \checkmark & ---        & ---        & Проприет. & Проприет. \\
Тестирование атак LLM      & \checkmark & ---        & ---        & ---        & ---        \\
MITRE ATT\MITRE ATT\&CK mappingCK~\cite{Hernandez2015} mapping      & \checkmark & Частично   & Частично   & \checkmark & \checkmark \\
Панель NIST AI RMF         & \checkmark & ---        & ---        & ---        & ---        \\
Постквантовая СОВ          & \checkmark & ---        & ---        & ---        & ---        \\
Открытый исходный код      & \checkmark & \checkmark & \checkmark & ---        & ---        \\
\bottomrule
\end{tabular}
\end{table}


%%=====================================================================
\section{Результаты валидации платформы}
\label{sec:platform_validation}
%%=====================================================================

\subsection{Производительность моделей}

Все 13+ моделей валидированы на бенчмарке CIC-IoT-2023 (Таблица~\ref{tab:platform_models}). Ансамбль SurrogateIDS достигает наивысшей точности (97,8\%) через семиветочную архитектуру. MambaShield обеспечивает конкурентную точность (95,9\%) при наименьшей задержке (0,5\,мс).

\subsection{Конвейер защиты LLM}

Четырёхуровневая система защиты~\cite{Anaedevha2026neurips} достигает макроусреднённого обнаружения 94,5\% по 517 сценариям атак. Абляция подтверждает, что каждый уровень обеспечивает независимый прирост: полный конвейер даёт восстановление на 41,6 п.\,п.\ точности извлечения RAG при 30\% контаминации корпуса.

\begin{figure}[H]
\centering
\includegraphics[width=0.92\textwidth]{figures/Adv_Eval_v10.png}
\caption{Страница состязательной оценки: (слева) панель конфигурации атак с выбором семейства атак, бюджета возмущений и целевой модели; (центр) деградация точности при FGSM, PGD и C\&W атаках с увеличением $\varepsilon$; (справа) сравнение стандартных и состязательно обученных моделей.}
\label{fig:screenshot_adversarial}
\end{figure}

\subsection{Пользовательское исследование}

Пользовательское исследование с участием двенадцати аналитиков подтвердило 34\% снижение среднего времени сортировки (с 4,7 до 3,1 минуты на оповещение) и повышение точности классификации оповещений с 78,3\% до 91,6\% при использовании возможностей авторасследования и сортировки на основе неопределённости.

\subsection{Валидационные наборы данных}

Три специально созданных набора PCAP-данных обеспечивают контролируемую эталонную базу для регрессионного тестирования:

\noindent\textbf{Набор А:} 60\,000 потоков, покрывающих все 34 класса атак со сбалансированным представлением классов.

\noindent\textbf{Набор Б:} 40\,000 потоков с акцентом на редких типах атак (Heartbleed, Infiltration, SSH-Patator).

\noindent\textbf{Набор В:} 50\,000 потоков с инъецированным дрейфом концепта на трёх темпоральных границах для валидации непрерывного обучения.

Постквантовые наборы данных получены из Kaggle (DOI: 10.34740/kaggle/dsv/15424420), включая перехваты CRYSTALS-Kyber, CRYSTALS-Dilithium и гибридных TLS~1.3 рукопожатий.


%%=====================================================================
\section{Вклад в открытый исходный код}
\label{sec:opensource}
%%=====================================================================

Полный исходный код платформы, все реализации моделей, конвейер защиты LLM и интерактивная панель управления архивированы в Zenodo (DOI: 10.5281/zenodo.19129512). Предобученные веса для архитектуры Multi-Agent PQC-IDS выпущены по адресу \texttt{github.com/rogerpanel/Multi-Agent-PQC-models}. Курированные постквантовые бенчмарк-наборы данных опубликованы на Kaggle (DOI: 10.34740/kaggle/dsv/15424420). Три воспроизводимых скрипта генерации PCAP обеспечивают детерминированное воссоздание оценочных наборов данных. Платформа лицензирована для академического и исследовательского использования, обеспечивая независимое воспроизведение всех представленных результатов.

\begin{figure}[H]
\centering
\includegraphics[width=0.96\textwidth]{figures/About_v10.png}
\caption{Страница <<О платформе>>, показывающая версию платформы, DOI-цитирование, информацию о разработчиках, технологический стек и ссылки на репозиторий GitHub, архив Zenodo и наборы данных Kaggle.}
\label{fig:screenshot_about}
\end{figure}


%%=====================================================================
\section{Итоги главы}
\label{sec:ch5_summary}
%%=====================================================================

В настоящей главе описана программная платформа RobustIDPS.ai, реализующая все методы, разработанные в данной диссертации, как единую развёртываемую систему. Платформа содержит 46 страниц в 9 функциональных группах, интегрирует 13+ моделей нейронных сетей, реализующих все семь методов (М1--М7) и их расширения, предоставляет четырёхуровневую систему защиты LLM с обнаружением 94,5\% по 517 сценариям атак с четырьмя коммерческими бэкендами LLM, включает модуль Multi-Agent PQC-IDS для классификации постквантового трафика и поддерживает интеграцию с существующими системами СОПВ (Snort, Suricata) через генерацию правил и потребление потоков оповещений. Платформа общедоступна как ПО с открытым исходным кодом (DOI: 10.5281/zenodo.19129512) и служит как операционным инструментом безопасности, так и воспроизводимым исследовательским инструментом. Пользовательское исследование подтвердило 34\% снижение времени сортировки аналитиками и улучшение точности классификации оповещений на 13,3 п.\,п. Платформа демонстрирует, что математические вклады, разработанные в Главах~\ref{ch:methods}--\ref{ch:experiments}, преобразуются в работающую реализацию, способную обрабатывать реальный сетевой трафик с операционными показателями пропускной способности и задержки.
